% !TeX program = xelatex
\documentclass[sectionlevel=book,status=draft]{noteformyself}

\title{Permanent Tag System Test Book}
\author{Tianle Yang}
\date{\today}
\covertitlefont{TeX Gyre Pagella}
\coversentence{A compact test document for published and draft entries.}
\version{1.2.0-dev}

% In a real notes workspace, this should be the single global registry shared by
% every repository. This local file is deliberately small enough for testing.
\LoadNoteTagRegistry{tag-registry.example.tex}

\notesetup{
  book={Tag System Test Book},
  book-tag=0,
  status=draft
}

\begin{document}

\maketitle
\tableofcontents

\chapter[
  status=published,
  tag=0AB,
  label=chap:published
]{The first tagged chapter}

\section[
  status=published,
  tag=0ABCD,
  label=sec:published
]{A published section}

\subsection[status=published,tag=0ABCDS1,label=subsec:published]{Published statements}

\begin{theorem}[
  title={Fundamental comparison},
  status=published,
  tag=0ABCDEF,
  label=thm:published-comparison
]
  If two objects satisfy the same universal property, they are uniquely
  isomorphic in the appropriate sense.
\end{theorem}

\begin{definition}[
  status=published,
  tag=0ABCDE1,
  label=def:published-object
]
  A published object is an entry whose semantic content and permanent tag have
  been recorded in the global registry.
\end{definition}

% The title key is genuinely optional, including for published theorems.
\begin{theorem}[
  status=published,
  tag=0ABCDE2,
  label=thm:untitled-published
]
  A permanent tag does not require the statement to have a title.
\end{theorem}

The local published structural references are \cref{chap:published},
\cref{sec:published}, and \cref{subsec:published}. The published statement
references are
\cref{thm:published-comparison}, \cref{def:published-object}, and
\cref{thm:untitled-published}.

The old identifier remains queryable as the frozen record
\NoteTagRecord{0ABCDE0}.

\chapter[
  status=draft,
  tag=0CD,
  label=chap:mixed-draft
]{A draft chapter}

% This is now the entry's current section. The proposition below was allocated
% under section tag 0ABEF and retains that old tag unchanged after the move.
\section[
  status=draft,
  tag=0CDEF,
  label=sec:mixed-draft
]{A mixed draft section}

\subsection[
  status=draft,
  tag=0CDEFS1,
  label=subsec:mixed-draft
]{Published and provisional entries}

This section remains draft because it contains both published and draft entries.

\begin{proposition}[
  status=published,
  tag=0ABEF01,
  label=prop:published-in-draft
]
  An individually published entry may remain inside a section that is still
  being drafted.
\end{proposition}

\begin{theorem}[label=thm:untagged-draft]
  This draft theorem has no provisional tag and therefore uses only its current
  structural location when referenced.
\end{theorem}

\begin{example}[
  tag=0ABEFZZ,
  label=eg:first-provisional
]
  A provisional draft tag is visible but is not reserved globally.
\end{example}

\begin{remark}[
  tag=0ABEFZZ,
  label=rmk:duplicate-provisional
]
  Repeating a provisional tag is permitted while the containing section remains
  draft. Publication validation will require a new unique registered tag.
\end{remark}

The draft structural references omit provisional tags:
\cref{chap:mixed-draft}, \cref{sec:mixed-draft}, and
\cref{subsec:mixed-draft}.

The mixed-section references are \cref{prop:published-in-draft},
\cref{thm:untagged-draft}, \cref{eg:first-provisional}, and
\cref{rmk:duplicate-provisional}.

% This negative test is disabled during an ordinary compilation. It can be
% enabled from the command line to confirm that strict publication fails.
\ifdefined\NFMTestMissingTag
  \section[
    status=published,
    tag=0CDIJ,
    label=sec:negative-missing-tag
  ]{Negative test: missing published tag}
  \begin{theorem}[status=published,label=thm:missing-published-tag]
    This entry must make a publication build fail.
  \end{theorem}
\fi

\ifdefined\NFMTestDraftInPublished
  \section[
    status=published,
    tag=0CDIJ,
    label=sec:negative-draft-entry
  ]{Negative test: draft entry in a published section}
  \begin{theorem}
    This draft entry must make a publication build fail.
  \end{theorem}
\fi

\end{document}
